% ============================================================
% 正文：Related Work and Background
% 作者：Kehui Chen, Jacky Keung, Xiaoxue Ma
% 主题：MCP 故障定位（fault localization）根因实证研究
% ============================================================

\section{Related Work and Background}\label{sec:related}

This section reviews the evolution of LLM tool integration, core MCP concepts, existing
fault-taxonomy research, and the gaps our work addresses.

\subsection{Evolution of LLM Tool Integration}

LLM-external tool integration has evolved through three phases to improve collaboration and
ecosystem compatibility \cite{liu2024toolace}.

\textbf{Early single-tool function calling.} Proposed by OpenAI's function-calling API, this
phase supports isolated tool invocations such as calculators and search engines
\cite{zhang2024large}. It cannot preserve cross-tool context, so task metadata must be
retransmitted for each call, making complex multi-step workflows unsupported
\cite{kim2024llm}\cite{wang2025function}.

\textbf{Second-stage agent orchestration frameworks.} Frameworks such as LangChain and
AutoGPT support chained tool calls through task planning and selection \cite{lin2025mao}.
Their proprietary protocols form isolated ecosystems and raise migration costs
\cite{haddad2017connecting}.

\textbf{Current-stage standardised middleware protocols.} As a standardised middleware, MCP
uses unified interfaces to break platform barriers \cite{anthropic_mcp_2024}\cite{hou2025model}.
It avoids vendor lock-in and enables cross-system collaboration for large-scale enterprise
deployment.

\subsection{The Model Context Protocol}

MCP is a standard middleware for LLM-tool collaboration. It resolves context fragmentation,
inconsistent interaction, and tight coupling to realise efficient, secure cross-system
cooperation. It consists of four core components (Host, Client, Server, and Resource
\cite{yang2025survey}), whose main features are:

\textbf{Unified context abstraction.} MCP adopts standard task-context formats for user
preferences, logs, and security policies, ensuring consistent data transmission and parsing.

\textbf{Decoupled layered architecture.} MCP decouples the LLM and tool layers so that each
can evolve independently, improving scalability and maintainability.

\textbf{Distributed dynamic context management.} MCP supports real-time context update,
sharing, and fine-grained access control to guarantee context security and consistency in
distributed systems.

\subsection{Current Research on MCP}

\textbf{MCP benchmarks.} A series of MCP benchmarks has been proposed for real-world LLM
evaluation, overcoming defects of early testing systems \cite{liu2025mcpeval}\cite{yang2025mcpsecbench}.

\textbf{MCP integration and performance.} Studies explore embedding MCP into existing systems
for performance improvement \cite{xu2025towards}\cite{li2025energyplus}\cite{da2025beyond}.

\textbf{Security and compliance.} MCP security is widely studied; researchers have classified
attacks and built defense frameworks \cite{hasan2025model}. Seventeen attack types and basic
protections such as access control and audit logging have been presented
\cite{yang2025mcpsecbench}\cite{mcpSecurity2025}\cite{fu2025mcp}. Most defenses remain passive.

\subsection{Fault Taxonomies and Fault-Oriented Benchmarks}

Recent work has begun to analyse failure modes in MCP deployments. Existing fault taxonomies
are mostly derived by mining GitHub issue reports and open-source repository discussions,
and identify several failure categories predominantly observed on the MCP Server side.
Security-oriented benchmarks such as MCPSecBench mainly construct adversarial attack
scenarios for vulnerability testing \cite{yang2025mcpsecbench}.

Nevertheless, existing fault-related work has three limitations. First, analyses based on
issue mining rely on static text reports rather than executable, end-to-end test cases, so
faults cannot be directly reproduced in controlled experiments. Second, existing taxonomies
place heavy emphasis on Server-side and security-oriented failures, giving limited attention
to faults originating from Host, Client, or Transport layers; ordinary operational runtime
faults and SDK observability behaviours remain underexplored. Third, existing studies
rarely investigate error-reporting phenomena such as manifestation-root-cause mismatch or
silent failure, which strongly affect the real-world debugging experience.

\subsubsection{Why a New Four-Layer Taxonomy? (Comparison with Existing Fault Taxonomies)}

To situate our work in the broader software-engineering literature, we build on the large-scale
empirical taxonomy of MCP faults recently proposed by Taraghi and co-workers
\cite{taraghi2026real}, who mined 3,282 bug-related GitHub issues and derived five high-level
fault categories (e.g.\ \texttt{Server/Tool Configuration} 31.74\%, \texttt{Server/Host
Configuration} 28.64\%, \texttt{Server Setting} 27.45\%), validated by a survey of 41
practitioners. Applied to our dataset, MCP faults would fall into a small number of these
categories: \emph{configuration-related} failures (triggered by erroneous
tool/parameter/credential configuration, e.g.\ \texttt{Server/Host Configuration},
\texttt{Server/Tool Configuration}) and \emph{general programming} failures (triggered by a
bug inside the tool's own execution logic, e.g.\ \texttt{General Programming}). Each reproduced
real bug in our dataset is mapped to such a category in Section~\ref{sec:setup}.

Crucially, Taraghi et al.\ identify \emph{diagnostic friction} as a key issue (their Findings
F4/F5 and insight A4) and recommend making host-side and protocol-boundary failures
\emph{observable}. However, their study is issue- and survey-based: it \emph{describes} which MCP
faults are hard to diagnose, but provides no controlled, reproducible mechanism to measure
\emph{why} diagnosis fails for specific fault types. Comparing their taxonomy with our four-layer
scheme highlights the gap we address. Taraghi's taxonomy classifies a fault by its \emph{trigger
type} and architectural \emph{boundary of resolution}, but it does not distinguish \emph{where
the fault originates} from \emph{where the fault is reported}. In an MCP pipeline, a single tool
call runs between the Host, Client, and Server components, and a fault may arise in any of them
or during transport between them; yet the SDK exposes the resulting error at only one layer.
Consequently, many MCP faults are triggered by a configuration or
programming error at one layer, but are \emph{reported} at the Server layer. The Taraghi
categories therefore describe the \emph{nature} of the defect but cannot express its
\emph{location}, which is exactly the information a developer needs when localising a fault.

Our four-layer taxonomy addresses this gap by classifying each fault by the architectural layer
at which it \emph{originates} (Host scheduling, Client message, Server tool, or Transport). This
makes two properties explicit that prior taxonomies obscure: (i) the \emph{cross-component}
nature of MCP faults, in which the reported layer differs from the originating layer, and (ii)
the \emph{manifestation-root-cause separation} that makes localisation uniquely difficult. The
two taxonomies are therefore \emph{complementary} rather than competing: Taraghi explains
\emph{what kind of bug} occurs and \emph{where it is resolved}, whereas our taxonomy explains
\emph{where it originates and how it is (mis)reported}. Moreover, whereas Taraghi stops at
recommending observability, our benchmark \emph{measures} it: we quantify the three channels
through which faults evade localisation and show that a runnable diagnostic-clue tool improves
it. We adopt the latter as our primary dimension because it directly supports fault
localisation and observability.

\subsection{Research Gaps and Motivation}

Existing MCP studies have clear limitations for enterprise deployment. First, most
benchmarks are laboratory-based and lack real enterprise data \cite{liu2025mcpeval}.
Second, integration research concentrates on technical embedding, ignoring practical
deployment challenges \cite{xu2025towards}\cite{li2025energyplus}. Third, security studies
focus on attack discovery and passive defence \cite{hasan2025model}. Fourth, empirical
studies relying purely on interviews or issue mining deliver subjective observations without
reproducible technical evidence to verify \emph{why} deployment pain points occur. Fifth,
there is a shortage of executable end-to-end fault datasets targeting general MCP runtime
failures, which hinders controlled experimental research on fault diagnosis and
observability.